\section{A Common Execution Contract}\label{sec:contract}
\subsection{Why the comparison object matters}
The historical corpus contains 9,360 configurations and two deterministic
verification passes per configuration. Its primary panel contains eleven
policies, with two additional style approximations. The historical engine
permits policy-dependent service modes: instantaneous prefill for some
policies, explicit prefill work for others, and differing enforcement of KV
limits. A policy's requested work can also differ from the engine's realized
progress. Consequently, those outcomes cannot be interpreted as a controlled
comparison of scheduling order. Table~\ref{tab:objects} distinguishes the
three evidence objects used in this paper.

\begin{table}[t]
\caption{Evidence objects and the conclusions they support. Historical
identifiers containing ``faithful'' are immutable names, not validation certificates.}
\label{tab:objects}\small
\begin{tabularx}{\linewidth}{@{}lXX@{}}\toprule
Object & Comparison or check & Interpretation \\\midrule
Historical LSSP & Policy-dependent service/resource modes; archived ranking arithmetic & Motivation and historical-rule counts \\
C01-FULL & Four adaptations; common \cm\ execution and information & Descriptive finite-corpus performance \\
STAT01 & Synthetic independent studies with known population truth & Scenario-specific inference calibration \\\bottomrule
\end{tabularx}
\end{table}

\subsection{Service, resources, and information}
\cm\ represents one abstract device. Each batch lasts 1\,ms and has a
32,768-token work budget. Each executed prompt token consumes one unit;
each selected decode request consumes one unit and emits one token.
All policies share these rules. A prefill batch emits no output, including
its final chunk: decoding can begin only in a later batch. Batch duration
is independent of its composition. Thus seconds and tokens per second
below are model quantities, not estimates for a particular GPU or LLM.

Each resident request reserves the complete 32,768-token context bound.
All frozen request lengths fit; the maximum prompt-plus-output length is
30,053. Capacities of 131,072 and 262,144 tokens therefore permit
\emph{four and eight resident requests}, respectively, even though the
configured sequence limit is 64. Actual used KV follows executed prompt
and output prefixes, rounded to 16-token blocks, and must remain within
the reservation. Completion releases the reservation. Waiting requests own
none. This conservative rule avoids output-length prediction and allocation
deadlock; it does not model incremental native paging or cache eviction.

There is no preemption, dropping, migration, prefix sharing, or recomputation.
Requests arriving during a batch become eligible at its end. Idle time
advances to the next arrival. The administrative horizon is the final
arrival plus 50\,s; a batch cannot finish beyond it. Unfinished requests
remain explicitly unfinished rather than contributing observed completion
latencies.

Policies see only arrived request identifiers, stable order, arrival time,
prompt length, acknowledged progress, residency, last-output/prefill-completion
anchor, and current time. Actual remaining output length, future arrivals,
source identity, and output-derived metric deadlines are hidden. The engine
either executes the exact proposed token spans or rejects the plan. The
shared contract therefore fixes both the resources and the online
information available to each adaptation.

\subsection{Adaptation identities and structural consequences}
Table~\ref{tab:policies} defines the four evaluated adaptations. Published
systems are lineage references; historical project implementations and
\cm\ adaptations are separate objects. The panel was targeted using the
historical comparison questions, not sampled to represent all schedulers.

\begin{table}[t]
\caption{Project-maintained adaptations. Artifact IDs append \texttt{\_c01}
to \texttt{whole\_fcfs}, \texttt{chunked\_fcfs}, \texttt{decode\_first}, and
\texttt{slo\_gated}. All use the shared contract above.}
\label{tab:policies}\small
\begin{tabularx}{\linewidth}{@{}p{.18\linewidth}p{.20\linewidth}X@{}}\toprule
Adaptation & Historical lineage & Executed ordering rule and principal restriction \\\midrule
Whole-FCFS & vLLM-inspired whole-prompt implementation & Admit whole prompts in arrival order when feasible; otherwise decode residents. Admitting prefills defers all decodes that batch. No native paged allocator. \\
Chunked-FCFS & vLLM-inspired chunked implementation & Resident work before waiting work in stable order; at most 512 prefill tokens per request per batch. \\
Decode-first & Sarathi-inspired implementation & Resident decodes, then resident prefills, then waiting prefills; 512-token cap. No original engine timing. \\
SLO-gated & SLAI-inspired implementation & Urgent decodes first; then resident prefills, waiting prefills by prompt length, and residual decodes. Fixed 0.495\,s gate; 512-token cap. No dynamic offset, heterogeneous classes, or full SLAI. \\\bottomrule
\end{tabularx}
\end{table}

Two consequences are essential to interpreting the results. At most eight
requests can execute, and each chunked request uses at most 512 tokens.
Their total work is therefore at most $8\times512=4,096$, below the shared
32,768-token budget. Every selected chunk and resident decode fits.
Chunked-FCFS and decode-first consequently make the same per-request
progress despite ordering their work differently. The complete saved
terminal records confirm this equivalence in all 1,440 matched conditions.

The same bound means that SLO-gated's residual decodes fit every batch.
After prefill completion, a resident's anchor advances with every output;
its elapsed time cannot reach the 0.495\,s urgency threshold in these
executions. The distinctive behavior is principally prompt-length admission
ordering with the frozen chunking rules. The results cannot be credited
to an active urgency gate. Whole-FCFS can still defer decoding for admitted
prefills; the work-bound equivalence does not apply to that rule.

\subsection{Conformance evidence}
The event referee reconstructs request states, token spans, timing,
completion, and resource use from the inputs and execution journal without
importing the engine or policy implementation. It checks that intended and
executed work agree and recomputes the metrics. Every campaign cell passed
this check; an independent analyzer also checked the complete ranking and
reversal reports. The final evidence review additionally checked all result
identities and terminal metrics and replayed 24 structurally selected new-cell
journals. These checks establish conformance to \cm, not GPU fidelity.
Figure~\ref{fig:contract} summarizes the evidence chain.

\begin{figure}[!htbp]
\centering
\begin{tikzpicture}[font=\small,node distance=5mm and 8mm,
 box/.style={draw,rounded corners=2pt,align=center,text width=.23\linewidth,minimum height=12mm},
 arr/.style={-{Latex[length=2mm]},thick}]
\node[box] (input) {Frozen windows\\and six settings};
\node[box,right=of input] (policy) {Four adaptations\\restricted online view};
\node[box,right=of policy] (engine) {Shared \cm\ engine\\time, work, KV};
\node[box,below=of input] (reports) {Absolute metrics\\matched comparisons};
\node[box,below=of policy] (referee) {Independent referee\\reconstructs execution};
\node[box,below=of engine] (journal) {Saved token journal\\and terminal records};
\draw[arr] (input)--(policy);\draw[arr] (policy)--(engine);
\draw[arr] (engine)--(journal);\draw[arr] (journal)--(referee);
\draw[arr] (referee)--(reports);\draw[arr] (input)--(referee);
\end{tikzpicture}
\caption{Comparison and verification under one model. The independent
analysis check recomputes aggregates, ranks, and reversals from the absolute
metrics. No stage in this chain establishes physical-system accuracy.}
\Description{Inputs pass through a restricted policy view and shared engine to a saved journal. An independent referee reconstructs execution before reporting metrics.}
\label{fig:contract}
\end{figure}
